\section{An observation contract from source records to a contact decision}
\label{exec13:observation-journey}

A servicing manager wants to know whether contacting an eligible borrower
reduces the probability of an adverse event during the following twelve
months. Before comparing contact with no contact, the analysis must establish
which borrower each record describes, what the adverse-event label means,
which borrowers could enter the observed sample, and which information was
available when contact was chosen. These requirements are related, but one
cannot stand in for another. A correctly dated row can belong to the wrong
borrower; a correctly linked record can carry a noisy label; and a table of
correct labels can represent a selected population. The proposed observation
contract makes these distinctions explicit before downstream causal analysis.

The immediate design uses generated laws and the existing local library and
command-line service. Its purpose is to establish what a minimal prototype
can preserve and which further operation is justified. It does not require
an expensive bank validation merely to explain a failure mechanism. The
development sequence is to specify the observation law, demonstrate the
consequence of violating it, implement a narrow supported path, and then ask
whether an unresolved bank-specific question could change the investment
decision~\citep{exec13_observation_design}.

\subsection{Keep the information clock separate from outcome ascertainment}

For a feature record, write event time $t_e$, first usable availability
$t_a$, and decision time $t_d$. The selected admission rule is
$t_e\leq t_a<t_d$. Equality of event and availability is permitted; equality
of availability and decision is not. A corrected value published after the
decision retains its later availability even when it refers to an earlier
event. Comparing timezone-aware instants avoids inventing a time ordering
from differently displayed local times. None of these checks establishes
that the supplied timestamps themselves are truthful.

An outcome has a different role. An event six months after contact can be
an eligible twelve-month outcome, although it must not appear among the
features available before contact. Its ascertainment date, censoring rule,
revision history and horizon belong in the outcome contract. Applying the
feature clock to every row would reject legitimate follow-up observations;
ignoring availability would admit future information. The proposed consumer
therefore declares the row's role before selecting an audit. The present
reference service implements the strict feature clock only and explicitly
withholds its alternative-clock operation~\citep{exec13_observation_primary}.

The proposed manifest also retains the source identity and owner declaration,
license and version, extraction and transformation history, origin and
household identifiers, candidate linkage maps, inclusion declarations,
label definition and adjudication evidence, and the counts of accepted,
rejected and unavailable records. A byte hash binds the content that was
used. It cannot establish who originally observed an event or whether the
contents are correct. Source custody, content integrity and scientific
measurement validity consequently have separate evidence fields and separate
failure reasons~\citep{exec13_observation_design}.

\subsection{Marginal inclusion probabilities do not determine precision}

Consider a fixed population with two origin values $y_1=20$ and $y_2=80$.
The equally weighted population mean is 50. Let $R_i$ indicate inclusion
and suppose both marginal inclusion probabilities are $\pi_i=1/2$. The
expansion estimator is
\[
 \widehat\mu=\frac12\sum_{i=1}^2\frac{R_i y_i}{\pi_i}
             =20R_1+80R_2.
\]
Its expectation is 50 under either independent inclusion or a shared
inclusion decision $R_1=R_2$. Under independent inclusion, its four equally
likely values are $0,20,80,100$, giving variance 1700. Under shared inclusion,
it equals zero or 100 with equal probability, giving variance 2500.
The marginal probabilities and expected sample size are identical in the
two experiments~\citep{exec13_observation_exact}.

More generally, for a fixed population of size $N$, known positive
inclusion probabilities and joint probabilities
$\pi_{ij}=E(R_iR_j)$, direct expansion gives
\begin{equation}
 \operatorname{Var}(\widehat\mu)
 =\frac{1}{N^2}\sum_{i,j}
   \frac{\pi_{ij}-\pi_i\pi_j}{\pi_i\pi_j}y_i y_j,
 \qquad \pi_{ii}=\pi_i.
 \label{exec13:observation-inclusion-variance}
\end{equation}
This is a design identity, obtained by expanding the variance of a weighted
sum. It is not an estimated variance from the observed sample, and it assumes
that the origin values and selection law have the stated meaning. In the
two-origin example the extra cross term is 800. Replacing the joint law
by independent selection because a file contains only the two marginal
probabilities would discard that term.

If the second origin is never included, populations $(20,80)$ and
$(20,120)$ induce the same observed law but have means 50 and 70. The
missing information is different from an imprecise estimate of a positive
inclusion probability. No finite replacement for $1/0$ identifies the unseen
value. Outcome bounds, a changed sampling design or additional information
can restrict the answer. If inclusion probabilities are fitted, their
estimation error and relation to the outcome contribution enter a new
uncertainty argument; Equation~\eqref{exec13:observation-inclusion-variance}
does not supply that argument automatically~\citep{exec13_observation_exact}.

Identity ambiguity has another geometry. Two received records with values
20 and 80 can be mapped bijectively to two origins in either order. The
total is 100 under both maps, but the second-minus-first origin contrast is
60 or $-60$. Independently drawing a match for each record could assign
both to one origin, which is outside this stipulated bijective class.
The consumer should retain the admissible joint maps, the target functional
and any evidence used to remove a map. These two values describe an authored
measurement contrast; they are not a treatment effect or a probability
distribution over all possible linkage mechanisms.

\subsection{Noisy audit labels can change the proposed action}

Let $Z$ be a true binary adverse event and $B$ its recorded label. Write
$p=P(Z=1)$, $s=P(B=1\mid Z=1)$ and $f=P(B=1\mid Z=0)$ for prevalence,
sensitivity and false-positive probability in the declared target population.
The law of total probability gives
\begin{equation}
 q=P(B=1)=f+(s-f)p.
 \label{exec13:observation-label-law}
\end{equation}
This relation is about measurement. It does not identify a contact effect.
If sensitivity or false positives vary with contact, borrower state or time,
the corresponding conditional laws must be specified and transported
appropriately. Correcting an aggregate prevalence cannot replace the
assignment and identification assumptions needed to compare interventions.

Two strictly positive authored joint laws illustrate the information gap.
The triples $(s,f,p)=(9/10,1/10,3/10)$ and $(1/2,1/10,3/5)$ both give
$q=17/50$. More records of $B$ alone cannot distinguish these two
measurement mechanisms. An audit may resolve the issue only to the extent
that it supplies suitable independent truth, known or justified inclusion,
and a relevant population. Calling the same noisy label a gold standard
does not create that information~\citep{exec13_observation_exact}.

Suppose, solely for a planning example, that contact prevents an event with
probability one half among borrowers who would otherwise experience it,
has no adverse event effect on the others, and costs USD~200. With an event
loss of USD~1000, the assumed expected incremental value is
\[
 V(p)=500p-200\quad\text{USD per eligible borrower}.
\]
The two equally compatible reported-label laws imply values $-50$ and
100. The sign changes even before sampling error is considered. The
one-half prevention mechanism is an additional stipulated causal model;
it is not a measured treatment effect inferred from $q$. The example
therefore motivates an experiment on the consequential missing information,
without asserting measured bank returns~\citep{exec13_observation_exact}.

\subsection{Propagate the fitted measurement law and preserve singular cases}

When $d=s-f\ne0$, the algebraic inverse is $p=(q-f)/d$. Suppose a
particular sampling and fitting design justifies a joint expansion of
$(\widehat q,\widehat s,\widehat f)$ with leading contributions
$(\psi_q,\psi_s,\psi_f)$. Differentiating the inverse gives
\begin{equation}
 \psi_p=\frac{\psi_q-p\psi_s-(1-p)\psi_f}{s-f}.
 \label{exec13:observation-label-score}
\end{equation}
For a declared common normalization, write the justified covariance of
those contributions as $\Omega$. Then the corresponding leading variance
is $g^\top\Omega g$, where
$g=(1,-p,-(1-p))^\top/d$. Different audit and main-sample sizes require
their proper scaling in the joint expansion. Household dependence, shared
label adjudication, reused fitting data and selection weights cannot be
omitted merely because each marginal estimate has a reported standard error.

An exact covariance illustration makes the cross terms visible. Take
independent unit-variance $X,Y$, and stipulate
$\psi_q=X$, $\psi_s=-X+Y$, $\psi_f=0$, $p=3/10$ and $d=4/5$.
Equation~\eqref{exec13:observation-label-score} then gives leading variance
$89/32$, whereas treating the fitted measurement law as fixed gives
$25/16$. These are authored influence contributions, separate from any
claim about the sampling law of the binary example. They show why a
covariance argument must concern the whole construction. They do not
establish that arbitrary audit estimators have this expansion
\citep{exec13_observation_exact}.

When $s=f$, the recorded label probability is $f$ for every $p$.
Near this regime the inverse and its derivative can be unstable. A more
faithful proposed interface carries a region $\mathcal C$ for the joint
observable and measurement parameters and reports
\begin{equation}
 \mathcal P(\mathcal C)=
 \{p\in[0,1]:\ \exists(q,s,f)\in\mathcal C,
                    \ q=f+(s-f)p\}.
 \label{exec13:observation-label-projection}
\end{equation}
If the true triple lies in $\mathcal C$, its true prevalence belongs to
this set, by substitution. Thus a genuinely justified joint-region coverage
statement transfers to the projected set. The projection cannot manufacture
that joint guarantee from an arbitrary collection of marginal intervals.
Nor does this argument establish a sharp identified set for a richer
observed law. The observation model, candidate region and target restriction
remain part of the claim.

For fixed $s=9/10$, $f=1/10$ and an expressly deterministic sensitivity
range $q\in[3/10,19/50]$, the compatible prevalence interval is
$[1/4,7/20]$. Equivalently,
$(4p-1)(20p-7)\leq0$. Mapping that whole interval through the illustrative
value function gives $[-75,-25]$ USD. At complete loss of label information,
the physical range $p\in[0,1]$ gives $[-200,300]$ USD. An empty compatibility
set instead indicates conflict among the supplied ranges and restrictions;
it is not evidence of zero loss. The generic local quadratic service can
represent finite intervals, empty sets, the whole line and disconnected
rays, but does not impose every target's physical range. The proposed
consumer must preserve the original result and apply justified target
restrictions explicitly~\citep{exec13_observation_exact,exec13_observation_primary}.

\subsection{What the local boundary comparison can establish}

The registered comparison fixes 26 authored cases across the library and
command-line routes, with an initial call and a restarted call through each.
Two further calls retain an existing key while changing the target horizon.
All 106 slots include refused, unsupported and failed requests as well as
computed results. Complete-response parity excludes only the request
correlation identifier. This tests a bounded local contract; it does not
estimate graph accuracy, sampling coverage or bank decision value
\citep{exec13_observation_design,exec13_observation_run}.

The cases deliberately expose the distinction between a completed diagnostic
and a validated observation. An empty row array can satisfy every declared
clock vacuously. Duplicate origins, false labels and an explicitly contradicted
identity statement can accompany valid timestamps. The present clock operation
checks their timing and retains its limited qualification; it does not repair
the other defects. Asking it jointly for clocks and identity does not silently
upgrade its scope. Measurement fitting, inclusion validation and noisy-label
truth remain separate unsupported requests~\citep{exec13_observation_primary,exec13_observation_run}.

Supplied-moment calculations reveal another boundary. An unverified measurement
statement can accompany a conditional numerical answer, while an explicitly
contradicted premise causes the selected estimator to refuse. The service
does not establish the truth of either declaration, automatically require
a preceding relevant audit, or connect that audit to a full causal analysis.
The prototype recommendation is therefore a consumer admission rule that
requires a nonempty audit of the right role and version, then either verifies
the remaining premises or propagates their uncertainty and limitations.
This admission rule is specified here; it is not claimed as an implemented
bank workflow~\citep{exec13_observation_design,exec13_observation_run}.

All 106 registered requests matched the frozen boundary expectations, and
all 26 cases preserved semantic responses across the two routes and restart.
The denominator contains 52 computed outputs, 20 unsupported requests,
four scientific refusals, 28 failed/unavailable results and two rejected
changed-horizon requests. A computed output includes diagnostics and generic
sets; this count is not the number of identified or useful causal estimates.
All negative dispositions remain in the comparison. Independent rational
evaluation checked 1,446 points against the supplied inequalities
\citep{exec13_observation_run}.

Total invocation wall time, including failed requests, replay and worker or
CLI launch, was 42.398 seconds in this local run. Active runner
time was 43.542 seconds, including fixture import and journaling.
These measurements exclude idle gaps between bounded segments, initial
runner startup, author review and all unimplemented work. Parent and child
memory records are separate process high-water readings, not simultaneous
process-tree memory. They do not establish a bank service level, a full
workflow price or a performance advantage over another product.

The journal preserved every earlier request across four execution segments
and a final aggregation-only resume. Restart reused the saved artifact
references and idempotency identities. It did not test process termination
at every possible publication point or multi-host recovery. No native
foreign-language engine, sibling repository, live model or authentic bank
owner participated in this screen~\citep{exec13_observation_run}.


\subsection{Choose the next information purchase by its decision consequence}

The first useful prototype combines the existing clock and artifact boundary
with one explicit observation model and a consumer that preserves unresolved
maps, selection and label uncertainty. A public dictionary can explain a
recorded field, and generated truth can test whether the interface loses
information. Neither proves that a bank's historical labels are correct.
Source navigation, proposed transforms and numerical operations should remain
separately accountable when an agent or another repository supplies them.
An actual sibling trial needs its own request, source/export identity,
response, failure/restart trace and owner disposition. A contract double
cannot substitute for that execution.

The complete cost account includes source acquisition and rights, extraction,
identity resolution, audit sampling and adjudication, fitting and covariance
work, data transfer, integration, scientific review, storage, failed requests,
monitoring and recovery. Local invocation time prices only a subset of that
work. Missing prices stay unknown. Before purchasing costly bank evidence,
the proposal asks whether the current feasible decision values differ enough
that improved identity, inclusion or label information could change the action.
If all admitted values already imply the same feasible action, more precise
measurement may have little incremental decision value. When signs differ,
the relevant study targets the uncertainty causing that difference, with
its cost compared against a justified expected benefit. The present generated
examples establish this rationale for a bounded prototype and an explicit
later experiment; they do not establish that any particular bank experiment
will be profitable~\citep{exec13_observation_design,exec13_observation_exact}.
\label{exec13:observation-journey-end}
